\chapter{Conclusion}\label{ch:conclusion}

This thesis implemented the \bine{} algorithms~\cite{desensi2025bine} inside \nccl{}, as four pieces: a tree-based broadcast and reduce (\cref{ch:tree}), a butterfly allgather and reduce-scatter (\cref{ch:butterfly}), the same allgather rebuilt on the \pat{} machinery (\cref{ch:pat}), and a tree-based allreduce that composes the first piece into a single fused kernel (\cref{ch:treeallreduce}).
Across the four, the experiments of \cref{sec:tree-results,sec:bfly-results,sec:pat-exp,sec:tar-results} draw a consistent picture of where the \bine{} structure pays for itself inside \nccl{} and where it does not.

\paragraph{Broadcast and reduce.}
The tree of \cref{ch:tree} wins in the band between the two regimes in which the ring is already good: below it, the ring with \texttt{LL} is latency-optimised and hard to beat (\cref{sec:tree-ll}); above it, with \texttt{SIMPLE} the ring becomes bandwidth-bound and its optimal bandwidth term takes over (\cref{sec:bg-ring-vs-tree}).
In between, from a few kB to a few MB and from $8$ ranks on, the logarithmic number of steps of the tree dominates, with speedups of $3.4$--$6.5\times$ over the ring at $128$ ranks (\cref{sec:tree-results}).

\paragraph{Allgather and reduce-scatter.}
The butterfly of \cref{ch:butterfly} only pays off in \emph{pos} mode, below the $n_cS$ crossover of \cref{eq:bfly-pos}, where a channel owns whole blocks and send and double send move a window as one contiguous, full-chunk transfer; in \emph{slice} mode, which is where every size above a few hundred kB falls, every strategy degrades to the calls of block-by-block while keeping its own extra cost, so none of send, permute or double send is ever worth it there (\cref{sec:bfly-res-modes}).
The measurements start at $1$\,kB, the smallest size exercised, where \bine{} already grows with $p$ by a constant instead of doubling like the ring (\cref{sec:bfly-res-ring}); nothing in the mechanism changes below that point, so the same gap, or a larger one, is expected for yet smaller blocks, only outside the range that was measured here.
Rebuilding the same schedule on the \pat{} machinery (\cref{ch:pat}) removes the dependence of the transfers on the data layout: all four strategies become equivalent, and block-by-block, the simplest, is the best one (\cref{sec:pat-res-modes}).
It does not, however, make the allgather faster in the configuration used throughout this thesis, four ranks per node: from $16$ ranks on \bine{}-\pat{} is $3$--$14\,\%$ slower than the butterfly kernels for large blocks and up to $2.2\times$ slower for small ones, and it is never the fastest algorithm once the ring may use \texttt{LL} (\cref{sec:pat-res-bfly,sec:pat-res-ring}).
It is competitive, and in one regime already ahead, only where the comparison is with \pat{} on its own terms, one rank per node, the only configuration the native \nccl{} \pat{} allgather supports (\cref{sec:pat-res-native}): there \bine{}-\pat{} is $1.55$--$1.74\times$ faster than the ring with \texttt{SIMPLE} and still ahead of the ring with \texttt{LL} at $16$\,kB from $8$ ranks on.
For small blocks the native \pat{} remains faster, thanks to a parallel factor and a fused receive-and-forward that \bine{}-\pat{} does not yet have, and the gap widens with $p$ over the tested range ($1.14\times$ at $4$ ranks to $1.47\times$ at $16$); this is a cost of the kernel, not of the schedule, so it is not expected to close with scale alone.
For blocks of $1$\,MB or more, instead, \bine{}-\pat{} is already $5$--$8\,\%$ \emph{faster} than the native \pat{} at the largest rank count tested, $16$, which the text attributes to the shorter average inter-node distance of the \bine{} schedule -- a structural property of the schedule, not an artefact of this range, so the advantage is expected to hold, and plausibly grow, at the larger rank counts this thesis did not have the allocation to reach.

\paragraph{Allreduce.}
The tree-based allreduce of \cref{ch:treeallreduce} roots a \bine{} tree differently on each channel and runs it as a single fused reduce-and-broadcast kernel.
The idea does not come from the reduce-scatter-plus-allgather composition of \cref{ch:butterfly}, whose windows halve and double at every step, but from the other allreduce of the \bine{} paper, the recursive doubling that exchanges the whole vector at every step with no such growth (Sec.~4.5 of~\cite{desensi2025bine}, \cref{sec:tar-motivation}): the multi-root tree can be read as that same pattern chopped down to $\min(p,\mathit{nchannels})$ simultaneous instances, one per channel, to match the number of channels \nccl{} actually runs, instead of one per rank (\cref{sec:tar-schedule}).
Of the four pieces, it has the best range: against the reduce-scatter-plus-allgather allreduce of the \bine{} butterfly, which does grow and shrink its windows, the tree is still the fastest algorithm at $1$\,MB on every scale and at $4$\,MB from $32$ ranks on, and the gap is largest exactly there: the butterfly allreduce is $1.3$--$3.4\times$ slower at $1$\,MB and $1.6$--$2.5\times$ slower at $4$\,MB, growing with the number of ranks (\cref{fig:tar-heatmap-simple}).
It also matches or beats the native \nccl{} tree up to $1$\,MB on almost every scale (\cref{sec:tar-results}).

\section{Future work}
Several directions follow directly from the limits above.

The standalone reduce and broadcast of \cref{ch:tree} process one child at a time in program order, a fresh single-peer \texttt{Primitives} object per level (\cref{sec:tree-collectives}), with the reduce in particular \emph{step-outer} (\cref{fig:bine-step-outer}).
This is a logical order, not a literal one: a send only blocks once the FIFO behind it is full, so consecutive levels already overlap in practice up to its depth, as \cref{sec:tree-collectives} notes for the broadcast.
A version with a single, wide-fan \texttt{Primitives} object per kernel, the same change that lets the fused allreduce of \cref{ch:treeallreduce} overlap its levels explicitly, was prototyped for both collectives; the preliminary runs showed an improvement only past the point where the ring already overtakes the tree (\cref{sec:tree-results}), i.e., outside the range where the tree is the better choice anyway.
Whether it helps inside that range is not established: a conclusive comparison would have to sweep the number of channels, the root, the message size, the number of ranks and the reduction operation against the existing kernels, a long sweep for the reduce in particular.
It is this evaluation, not the kernel change, that is left for future work.

The allgather rebuilt on the \pat{} machinery (\cref{ch:pat}) is also left short of what the machinery supports: it runs with a parallel factor of $1$, a single group of workers executing every piece in order, and without the fused receive-and-forward that keeps a piece in its slot and forwards it from there (\cref{sec:pat-kernel}).
Both are already present in the native \pat{} allgather and are what \cref{sec:pat-res-bfly,sec:pat-res-native} point to as the cause of the remaining gap at small blocks; nothing in the \bine{} schedule ties its operations to a single worker group, so raising the parallel factor and fusing the forward are natural next steps, not a redesign.

A further direction follows from the comparisons throughout this thesis, which are against \nccl{}'s own ring and double binary tree, never against a plain, non-\bine{} binomial tree under the same kernels and the same system.
Implementing the classic distance-halving binomial tree (\cref{sec:bg-bine}) as a fourth \nccl{} algorithm, alongside \bine{}, would isolate the gain that comes specifically from \bine{}'s shorter average distance between communicating ranks from the gain of using a pipelined tree at all, which the comparisons of this thesis do not separate.

The tree-based allreduce also leaves two things on the table.
The multi-root tree uses a single \bine{} tree per channel, whose root has $\log_2p$ children and therefore a bandwidth term that grows with $p$; combining two complementary trees, as \nccl{} does with its double binary tree (\cref{sec:bg-double-tree}), would remove this dependence for large vectors, which \cref{sec:tar-motivation} already identifies as the chapter's main limitation, and is left there as future work.
The second tree is not the obstacle, since the construction of \cref{sec:bg-bine} already gives one for any root; the obstacle is the load-balancing step that makes \nccl{}'s own double binary tree work, which books only half the inter-node bandwidth at the two boundary GPUs because every rank is a leaf with a fixed, small arity in one tree and an internal node in the other (\cref{sec:nccl-graph}).
A \bine{} tree does not have a fixed arity, so finding an equally effective balancing rule for it, one that holds across channels, roots and ranks rather than for a single fixed pair of trees, is the biggest hurdle to a \bine{} double tree, not the tree itself.

And the root of each channel is currently chosen once, by the bit-reversal permutation of \cref{sec:tar-schedule}, which only spreads the $\min(p,\mathit{nchannels})$ roots apart; a schedule that also balances the fan-out and the link load across channels, rather than only their position, could close part of the remaining gap with the native tree on the scales where it still wins.
The choice of graph is itself worth revisiting: among the alternatives tried, the default ring graph of \cref{sec:nccl-graph} proved more consistent across topologies and allocations, and its permutation minimises the hop-to-hop distance around the whole cycle, wrap-around edge included, which only approximates minimising the modular distance that \bine{} actually cares about.
A search built around \bine{}'s own edges, the specific pairs its tree and butterfly connect at each step, rather than the ring's adjacent pairs, might do better than this coincidence, though whether it would generalise as consistently as the ring graph does across topologies and allocations is itself an open question.

Last, the \emph{pos}/\emph{slice} split of the butterfly (\cref{sec:bfly-pos}) exists only because a window doubles, or halves, in size at every step of the allgather and the reduce-scatter, while the partition of \nccl{} is fixed across channels before the kernel starts (\cref{sec:nccl-cbd}).
\Cref{ch:pat} already tried the kernel-side answer, aggregating the steps on the \pat{} machinery instead of repartitioning, and it removes the dependence on the data layout without making the allgather faster (\cref{sec:pat-res-bfly}).
The partition itself is therefore the lever left: a channel repartitioner that reassigns the work between channels at run time, as the window grows or shrinks from one step to the next, would let the contiguous, full-chunk transfers of \emph{pos} mode apply past the few hundred kB where they give out today, at the cost of changes to the continuous byte distribution of \nccl{} (\cref{sec:nccl-cbd}) that go well beyond the kernels this thesis touched.
